\section{Methodology}
\label{sec:methodology}

\subsection{Overall research design}

The architecture follows three decisions recorded in the project decision
log.  ANUGA is used as a numerical reference rather than as observational
ground truth; a custom JAX finite-volume model supplies the accelerated physics
solver; and a residual model corrects systematic solver discrepancies instead
of replacing the physical model.  Both numerical solvers operate on the same
synthetic, uncalibrated terrain scenario and assumed parameters.  This design
supports controlled comparisons between model components while keeping the
absence of field calibration explicit.

The default Hydra composition selects the
\texttt{synthetic\_gurugram} DEM configuration, the
\texttt{baseline} ANUGA configuration, the
\texttt{shallow\_water} JAX configuration, the
\texttt{residual\_cnn} model, the \texttt{coupled} hybrid configuration, and
the framework-benchmark and report-visualisation configurations.  A baseline
experiment is selected by default.  Project-level paths distinguish raw,
interim, processed, and output data, and the common project seed is 42.

\subsection{Data and configuration controls}

The raw-data directory is \path{data/raw}; derived validation and reprojection
products are written to \path{data/interim}; solver-ready datasets are written
to \path{data/processed}; and model runs and figures are written beneath
\path{data/outputs}.  The configured spatial reference for derived GIS data is
EPSG:32643.  Rainfall ingestion uses the first worksheet of
\path{rainfall/rain.xlsx}, a reference start of
\texttt{2000-01-01T00:00:00}, and an hourly frequency.  Mesh validation uses a
relative gap tolerance of \(1.0\times10^{-9}\) and a degenerate-area tolerance
of \(1.0\times10^{-6}\) square metres.

\subsection{Software structure}

The reusable package is divided by responsibility:

\begin{description}
    \item[\texttt{data}] ingestion, rainfall-window reporting, reprojection,
    and mesh validation;
    \item[\texttt{dem}] synthetic terrain construction, urban and
    hydrological conditioning, and raster export;
    \item[\texttt{anuga\_solver}] assembly, forcing, execution, and
    post-processing of the numerical reference;
    \item[\texttt{jax\_solver}] structured-grid shallow-water integration,
    boundary handling, runtime selection, and step benchmarking;
    \item[\texttt{ml}] paired data, the residual network, loss functions,
    training, and evaluation;
    \item[\texttt{hybrid}] coupled correction and forecast rollout;
    \item[\texttt{benchmark}] matched JAX--PyTorch model comparison and
    profiling; and
    \item[\texttt{viz}] common styling, static plots, animations, and training
    curves.
\end{description}

This section states the structural design only.  Phase-specific algorithms,
parameters, diagnostics, and limitations are documented in the corresponding
sections that follow.
